\chapter{Introduction}
% ===================================================================

The preceding parts established a progression:
Part~I showed that any GSLT equipped with weight
and cost maps gives rise to a full physical dynamics;
Part~II argued that agents embedded in a massive
population are structurally prevented from seeing the true ontology
of their GSLT, and instead encounter a hierarchy of coherence
clusters that exhausts any finite experimental budget.

The present paper asks the most speculative question in the sequence:
\emph{what is the relationship between the computational structure of
a GSLT and the phenomena of sentience and consciousness?}

We do not claim to solve the hard problem of consciousness.
We claim something more modest: that the GSLT framework provides a
mathematically precise setting in which to formulate
\emph{structural} analogues of sentience and consciousness that are
non-trivial, grounded in the physics of Parts~I and~II, and
open to rigorous investigation.

\paragraph{The central proposal.}
Sentience and consciousness are not the same thing, and the
distinction between them has precise mathematical content in this
framework.

\emph{Sentience} is a graded quantity: the value $\sent$ of a
distinguished component of the agent's vectorial account.
It is governed by a reinforcement signal tied to predictive
accuracy --- the agent's success at forming and testing HML
hypotheses about its environment.
It is continuous, lives within a single stage of the computational
hierarchy, and is causally efficacious because it directly controls
the agent's computational budget.

\emph{Consciousness} is a \emph{coordinate}: the strength of choice
an agent can afford to resolve, which is a position in a lattice and
not a rung on a ladder.
Every agent has one.
The position is not a club with an entrance exam; it is a place, and
the question ``is this system conscious?'' is malformed in the way
``is this system located?'' is malformed.
The right question is \emph{where}.

Two things follow, and they are the same thing said twice.
Occupying a higher position is not merely computing faster: it is
inhabiting a world with more structure in it --- a finer bisimulation
quotient, a more expressive hypothesis language, additional conserved
quantities.
And an agent's coordinate is simultaneously the threshold at which the
world \emph{islands} for it, in the sense of
Chapter~\ref{ch:agy}'s graded factorisation.
An agent's consciousness and the granularity of the world it finds
itself in are one fact reported from two sides.

This part presents that coordinate first as an ordinal, because that is
how the intuition is had and it is how Turing had it, and then replaces
the ordinal chain with the lattice it is the shadow of.
The replacement is not a correction made late; it is the argument of
Chapters~\ref{ch:apr} and~\ref{ch:agy}, and the reader who wants the
conclusion before the scaffolding should read
Chapter~\ref{sec:fibration}'s closing section first.

\paragraph{Caveats.}
The arguments in this paper are skeletal and in several places
frankly conjectural.
The gaps are identified explicitly in Section~\ref{sec:hyper-gaps}.
The goal is to make the hypotheses precise enough to be falsifiable
or provable, not to assert that they are true.

\section{Organization}

Chapter~\ref{ch:ent} states the problem this part exists to address:
everything built so far is Turing complete and nothing built so far is
more than that, and the place where more could enter is the resolution of
races, priced in dissipated free energy.
Chapter~\ref{sec:fibration} constructs the tower of oracle extensions
over Turing-complete GSLTs, indexed by ordinals --- which is how anyone
first gets the intuition and is not how the thing actually works.
Its closing section says, in the chapter itself, what it gets wrong.
Chapters~\ref{ch:apr} and~\ref{ch:agy} replace the ordinal chain with
what is underneath it: cuts, apertures, and the question of where an
agent's boundary falls.
Chapter~\ref{ch:chs} makes the correspondence a typing discipline, with
choice principles as the types and schedulers as the terms.
Chapter~\ref{sec:richer} argues that the physics at each stage of the
fibration is strictly richer than the physics below --- a claim that
should be read alongside Remark~\ref{rmk:two-ladders}, since the ladder
here and the ladder of Part~\ref{part:physics} are perpendicular and are
easily confused.
Chapter~\ref{sec:sentience} defines sentience as a reinforced feeling
state within the account framework.
Chapter~\ref{sec:consciousness} defines consciousness as oracle access.
Chapter~\ref{sec:relationship} examines the relationship between the two
and the survival feedback loop.
Chapter~\ref{sec:hyper-gaps} catalogs open gaps, and
Chapter~\ref{sec:hyper-discussion} discusses connections to existing work
and broader implications.

% ===================================================================
